\documentclass[11pt]{article}
\usepackage[margin=1in]{geometry}
\usepackage{amsmath,amssymb,booktabs,array}
\usepackage[hidelinks]{hyperref}

\title{Miscellaneous Checks: Event Studies, Paired Movers, and Kline's Edge Model}
\author{}
\date{\today}

\newcommand{\E}{\mathbb{E}}
\newcommand{\Var}{\operatorname{Var}}
\newcommand{\Cov}{\operatorname{Cov}}

\begin{document}
\maketitle

\section*{Source Convention}

Page citations to Sorkin--Warwar refer to the PDF pages in
\texttt{sorkin warwar 2026 sw\_logadditivity.pdf}. Page citations to Kline refer
to the PDF pages in \texttt{kline\_2024\_review\_akm\_firm\_wage\_effects\_w33084.pdf}.
When I reuse the prior project write-up, I cite the prior result PDF
\texttt{7\_15\_Meeting (1).pdf} and \texttt{shimer\_smith\_writeup.tex}.
Displays or claims not directly sourced are labeled as derived examples,
derived algebra, or derived mappings.

\section{Notation}

\begin{center}
\begin{tabular}{>{\raggedright\arraybackslash}p{0.24\linewidth}
                >{\raggedright\arraybackslash}p{0.68\linewidth}}
\toprule
Symbol & Meaning \\
\midrule
$Y_{it},Y_{it}(j)$ & Observed wage and potential wage for worker $i$ at time $t$ if assigned to firm $j$. Source for Kline potential-outcome notation: Kline, pp. 29--31. \\
$D_{it}$ & Firm employing worker $i$ at time $t$ in Kline's causal notation. Source: Kline, p. 30. \\
$\alpha_i,\psi_j,X_{it}'\beta,\varepsilon_{it},\mu_{ij}$ & Worker effect, firm effect, covariate component, AKM error, and illustrative time-invariant worker-firm match effect. Source for AKM components: Kline, p. 15. The match-effect extension is derived for the separability discussion. \\
$\tau_t,\varepsilon_{it}(j)$ & Illustrative common time component and transitory potential-wage shock used in the derived parallel-trends example. \\
$A,B,a,b$ & Origin and destination firms and two paired movers in the derived paired-mover algebra. Source for paired-mover design: SW abstract, p. 1. \\
$G_A,G_B,\beta_{\mathrm{paired}}$ & Origin wage gap, destination wage gap, and paired-mover forecast coefficient. Derived notation for the paired-mover separability calculation. \\
$h_i,p_j,m_{ij}$ & Worker, firm, and match effects in SW's partial-equilibrium search model. Source: SW, p. 27. \\
$EE,EUE,\lambda_1$ & Employer-to-employer moves, employer-unemployment-employer moves, and SW's on-the-job offer-arrival parameter. Source: SW, pp. 27, 31--33. \\
$X,Y,g,\beta_{\mathrm{pooled}}$ & Generic regressor, outcome, group, and pooled OLS coefficient in the derived pooling identity. \\
$R,E,B,\Delta,u,\hat\Delta,\tilde\Delta$ & Kline adjusted wage changes, worker-by-edge matrix, firm-by-edge incidence matrix, unrestricted directed edge effects, edge-model error, estimated edge effects, and AKM-predicted edge effects. Source: Kline, pp. 18--19, 24--27. \\
$c,C,\dot\psi,\dot\eta$ & Kline cycle vector, matrix collecting fundamental cycles, projected firm-effect coefficients, and cycle-effect coefficients. Source: Kline, pp. 18--20. \\
$W,H,M,h_{\ell\ell},n_\ell,\ell$ & Kline mover-weight matrix, projection matrix, residual-maker matrix, edge leverage, number of movers on edge $\ell$, and edge index. Source: Kline, pp. 24--27. \\
$Y_{ij},X_{ij}'\delta,u_i,v_j,\Lambda$ & Project low-rank wage-interface notation used in the nesting comparison. Source: prior result PDF, pp. 18--21; \texttt{shimer\_smith\_writeup.tex}, Sec. 6. \\
$m_{ij},\mu_{jk}$ & Derived static mean wage surface and edge-specific mean worker factor, $\mu_{jk}=\E[u_i\mid i:j\to k]$, used to map low-rank surfaces into Kline edge effects. \\
\bottomrule
\end{tabular}
\end{center}

\section{Sorkin--Warwar: Why the Event Study Can Have Low Power}

The intuition in the to-do list is right, but it needs one sharpening. The standard
event-study diagnostic is mainly a diagnostic for pre-trends and realized mover
wage changes lining up with firm-effect changes. It is not a direct test of additive
separability. A time-invariant match effect can sit in the level of wages and pass
a parallel-trends diagnostic cleanly.

\subsection{Parallel Trends Is Not Separability}

Derived illustrative model: consider the potential wage level at firm $j$:
\[
Y_{it}(j)=\alpha_i+\psi_j+\mu_{ij}+\tau_t+\varepsilon_{it}(j),
\]
where $\mu_{ij}$ is a time-invariant worker-firm match effect. After removing
common time effects, the origin-firm potential trend for a worker moving from
$j$ to $k$ is
\[
Y_{i2}(j)-Y_{i1}(j)
=
\varepsilon_{i2}(j)-\varepsilon_{i1}(j).
\]
If the move is not selected on this transitory shock, then
\[
\E[Y_{i2}(j)-Y_{i1}(j)\mid D_{i1}=j,D_{i2}=k]=0.
\]
This is exactly Kline's parallel-trends restriction: among workers switching
between a pair of firms, average potential wages at the origin firm would not
have changed between periods. Source: Kline, Assumption 2, p. 30.

Derived comparison: additive separability requires something stronger:
\[
\mu_{ij}=0
\quad\text{or at least}\quad
\E[\mu_{ik}-\mu_{ij}\mid i:j\to k]=0
\text{ in the relevant contrasts.}
\]
The realized wage change for a mover is
\[
Y_{i2}(k)-Y_{i1}(j)
=
\psi_k-\psi_j+\mu_{ik}-\mu_{ij}
+\varepsilon_{i2}(k)-\varepsilon_{i1}(j).
\]
So a time-invariant match effect does not create a pre-trend by itself. It creates
heterogeneous treatment effects and possible intransitive edge effects. That is a
different restriction from parallel trends.

This is precisely Kline's point in Section 3.3: edge wage changes can have a
causal interpretation under exclusion, parallel trends, and stationarity even when
additivity fails. Source: Kline, pp. 29--31. He explicitly notes that one does not
need additive separability for firm switches to identify average treatment effects
on movers. Source: Kline, Proposition 1, p. 31.

\subsection{Why Paired Movers Hit Separability Directly}

Derived paired-mover algebra: take two workers $a$ and $b$ who both move from firm $A$ to firm $B$. Their
origin wage gap is
\[
G_A
=
Y_{aA}-Y_{bA}
=
(\alpha_a-\alpha_b)+(\mu_{aA}-\mu_{bA}),
\]
and their destination wage gap is
\[
G_B
=
Y_{aB}-Y_{bB}
=
(\alpha_a-\alpha_b)+(\mu_{aB}-\mu_{bB}).
\]
The firm effect cancels because both workers traverse the same $A\to B$ pair.
Under additive separability, $\mu_{ij}=0$, so
\[
G_B=G_A.
\]
With match effects, $G_A$ and $G_B$ need not be equal. Derived special case: if match effects are
independent across firms and independent of worker effects, then
\[
\beta_{\text{paired}}
=
\frac{\Cov(G_B,G_A)}{\Var(G_A)}
=
\frac{\Var(\alpha_a-\alpha_b)}
{\Var(\alpha_a-\alpha_b)+\Var(\mu_{aA}-\mu_{bA})}
<1.
\]
This is the clean mathematical reason paired movers test separability more
directly than the event study. The paired-mover comparison removes the common
firm move and asks whether worker wage gaps are portable across the same
origin-destination transition.

SW's headline empirical design is exactly this: for two workers both moving from
firm $A$ to firm $B$, the pay gap at the origin should predict the gap at the
destination with coefficient one under additive worker and firm effects. Source:
SW abstract, p. 1. They estimate coefficients below one in Brazil and Italy, which
they interpret as evidence of match effects. Source: SW abstract, p. 1.

\subsection{The Low-Power Mechanism in SW}

SW's model gives workers payoff
\[
h_i+p_j+m_{ij},
\]
where $h_i$ is a worker effect, $p_j$ a firm effect, and $m_{ij}$ a match effect.
Workers observe both firm and match components when deciding whether to accept
offers, so the model violates both additive separability and exogenous mobility.
Source: SW, p. 27.

Even so, the pooled event study can pass. In their simulated model, the full-sample
event-study coefficient is $0.99$ even though $41\%$ of moves are employer-to-employer
moves. Split by move type, the coefficient is $0.57$ for EE moves and $0.99$ for
EUE moves. Source: SW, pp. 31--33. So the event study detects the failure in the
EE subsample, but the pooled sample hides it.

Derived pooling identity: if the sample is split into groups $g$, the pooled
OLS slope of $\Delta Y$ on $\Delta\psi$ is not just a sample-size-weighted average
of within-group slopes. It also includes between-group covariance:
\[
\beta_{\text{pooled}}
=
\frac{\sum_g \Pr(g)\Cov(X,Y\mid g)
+\Cov(\E[X\mid g],\E[Y\mid g])}
{\sum_g \Pr(g)\Var(X\mid g)+\Var(\E[X\mid g])},
\qquad X=\Delta\psi,\;Y=\Delta Y.
\]
The between-group slope can push the pooled coefficient back toward one even when
within-group slopes show failure. SW explicitly describe this mechanism: pooled
coefficients combine relative sample sizes, within-group variances, and between-group
differences in sample means. Source: SW, pp. 32--33. In their simulations, increasing
the EE share lowers both subsample coefficients, but the pooled coefficient remains
above $0.90$ and can exceed either subsample coefficient because the between-sample
slope is above one. Source: SW, pp. 32--34.

\subsection{Bottom Line on the Friday Intuition}

Derived bottom line from Kline and SW: the Friday intuition is correct:
\[
\text{event study} \approx \text{parallel-trends / realized mover diagnostic},
\]
while
\[
\text{paired movers} \approx \text{direct portability / separability diagnostic}.
\]
A time-invariant match effect can pass the event study because it does not create
origin-firm pre-trends. It fails paired movers because the worker wage gap need not
be preserved from $A$ to $B$. SW's additional contribution is showing that even when
endogenous mobility creates event-study failures in a subsample, pooling can make the
full event-study coefficient look reassuring. Source: SW, pp. 31--35; prior result
PDF, Sec. 4 and Sec. 7.7.

\section{Kline Section 3: What the Content Is}

Kline Section 3 rewrites AKM as a restriction on directed mover edge effects. Start
from
\[
Y_{it}=\alpha_i+\psi_{j(i,t)}+X_{it}'\beta+\varepsilon_{it}.
\]
Source: Kline, p. 15. For two-period movers, define adjusted wage changes
\[
R=Y_2-Y_1-(X_2-X_1)'\beta.
\]
Let $E$ be the worker-by-edge matrix of directed origin-destination indicators,
and let $B$ be the firm-by-edge incidence matrix. Kline shows that AKM implies
\[
R=EB'\psi+\varepsilon.
\]
Source: Kline, pp. 18--19.

The unrestricted edge-effect model is
\[
R=E\Delta+u,
\]
where $\Delta_{jk}$ is the average adjusted wage change for movers from firm $j$
to firm $k$. Source: Kline, p. 19. AKM imposes
\[
\Delta=B'\psi.
\]
This reduces potentially $|E|$ directed edge effects to $J-1$ independent firm
effects. Source: Kline, p. 19.

\subsection{The Cycle Content}

For any cycle vector $c$ in the mobility graph, $Bc=0$. Therefore AKM implies
\[
c'\Delta=c'B'\psi=0.
\]
Source: Kline, pp. 18--19. Conversely, within a connected set, these cycle
restrictions exhaust the restrictions that AKM imposes on edge effects. Kline
writes
\[
\Delta=B'\dot\psi+C\dot\eta,
\]
where $C\dot\eta$ is the cycle component, and AKM is the restriction
$\dot\eta=0$. Source: Kline, p. 20.

This is the actual content of Section 3: firm effects are not primitive magic
objects. They are a low-dimensional potential function for directed mover wage
changes. They are valid as global firm wage levels only when edge effects are
path-independent around cycles.

\subsection{The Noise-Adjusted Goodness-of-Fit Content}

Kline also asks how well AKM approximates unrestricted edge effects in Veneto.
The key object is the residual between estimated unrestricted edge effects and
AKM-predicted edge effects:
\[
\hat\Delta-\tilde\Delta=M\hat\Delta,
\qquad M=I-H.
\]
For the mover-weighted residual sum of squares, Kline decomposes
\[
\E_u[\hat\Delta'M'WM\hat\Delta]
=
\Delta'M'WM\Delta
+\operatorname{trace}\!\left(M'WMV_u[\hat\Delta]\right).
\]
Source: Kline, p. 27. The first term is model error; the second is sampling noise.
Under exact AKM, the first term vanishes but the second remains.

Empirically, among non-bridge edges with at least two movers, the residual sum of
squares is $360.51$, while expected noise under AKM is $207.58$. The difference
$152.92$ is the estimated squared approximation error. Relative to the
noise-adjusted edge-effect sum of squares $978.67$, AKM captures about $84\%$ of
true edge-effect variation; the implied correlation is $0.92$. Source: Kline,
pp. 27--28. Across edge samples, Kline reports noise-adjusted $R^2$ values roughly
between $72\%$ and $88\%$, and about $81\%$ when bridges are included. Source:
Kline, pp. 28--29.

\subsection{The Causal Content}

Kline's causal section says that edge effects can be causal under three conditions:
exclusion, parallel trends, and stationarity. Source: Kline, pp. 29--31. Proposition
1 states that, under those assumptions,
\[
\E[Y_{i2}-Y_{i1}\mid D_{i1}=j,D_{i2}=k]=\Delta_{jk},
\]
the average treatment effect on movers along edge $j\to k$. Source: Kline, p. 31.
Crucially, this does not require additive separability. Kline explicitly emphasizes
that edge effects can be causal even when firm effects are heterogeneous, but those
edge effects may be intransitive and therefore may not aggregate to a single global
firm ranking. Source: Kline, pp. 30--31.

\section{Is Kline's General Edge Model Nested in Our Low-Rank Factor Model?}

Short answer:
\[
\text{AKM is nested; unrestricted Kline edge effects are not.}
\]

\subsection{AKM Is Nested}

The additive AKM wage surface is the rank-zero case of the project interface
\[
Y_{ij}=X_{ij}'\delta+\alpha_i+\psi_j+u_i^\top\Lambda v_j+\varepsilon_{ij},
\qquad
\Lambda=0.
\]
This was already in the project note and prior result PDF. Source: prior result
PDF, pp. 18--21; \texttt{shimer\_smith\_writeup.tex}, Sec. 6.

\subsection{A Low-Rank Wage Surface Implies a Restricted Edge Model}

Derived mapping from the project interface to Kline edges: suppose the true static wage surface is
\[
m_{ij}=\alpha_i+\psi_j+u_i^\top\Lambda v_j.
\]
For a worker moving from firm $j$ to firm $k$, the wage change is
\[
m_{ik}-m_{ij}
=
\psi_k-\psi_j+u_i^\top\Lambda(v_k-v_j).
\]
The mean edge effect is therefore
\[
\Delta_{jk}
=
\psi_k-\psi_j+\mu_{jk}^\top\Lambda(v_k-v_j),
\qquad
\mu_{jk}\equiv \E[u_i\mid i:j\to k].
\tag{LR-edge}
\]
This is the bridge from our low-rank wage surface to Kline's edge language. Around
a directed firm cycle $C$, the additive firm effects telescope:
\[
\sum_{(j,k)\in C}\Delta_{jk}
=
\sum_{(j,k)\in C}\mu_{jk}^\top\Lambda(v_k-v_j).
\]
So nonzero cycle components can arise from two ingredients together: nonzero
interaction $\Lambda$ and edge-specific mover composition $\mu_{jk}$.

\subsection{But Kline's Unrestricted Edge Model Is More General}

Kline's unrestricted model
\[
R=E\Delta+u
\]
places a free parameter on each observed directed firm-to-firm edge. It is not a
static worker-firm wage surface. It is a reduced-form model of average wage changes
for selected movers along origin-destination edges.

Therefore it is not nested in $(\star)$ unless extra restrictions happen to hold.
For a Kline edge vector $\Delta$ to come from a low-rank wage surface, there must
exist firm effects $\psi_j$, firm factors $v_j$, interaction matrix $\Lambda$, and
edge-specific mover composition moments $\mu_{jk}$ such that (LR-edge) holds for
every observed edge. An arbitrary $\Delta$ need not satisfy this.

Equivalently:
\[
\text{low-rank wage model}
\quad\Rightarrow\quad
\text{structured edge effects},
\]
but
\[
\text{unrestricted edge effects}
\quad\not\Rightarrow\quad
\text{low-rank wage model}.
\]

Kline's $\Delta_{jk}$ is therefore best treated as a diagnostic/statistical interface.
It can absorb the implications of additivity, low-rank interactions, match effects,
selection, or other mechanisms for mover wage changes. But by itself it does not
identify which mechanism generated those edge effects. This matches the prior
result PDF's conclusion that Kline's edge effects are diagnostic/statistical-interface
objects, not structural models. Source: prior result PDF, Secs. 7.4--7.7 and 8.3.

\section{Clean Takeaways}

\begin{enumerate}
\item The event study is mainly a parallel-trends / realized-mover diagnostic. It can
pass with time-invariant match effects because those match effects sit in wage levels
and treatment-effect heterogeneity, not necessarily in pre-trends.

\item Paired movers test separability more directly. For two workers making the same
$A\to B$ move, firm effects cancel, so the test asks whether wage gaps are portable
from origin to destination.

\item SW's additional low-power mechanism is pooling. Within-subsample event-study
coefficients can reveal problems, while between-sample mean differences push the
pooled coefficient back toward one.

\item Kline Section 3 says AKM is a path-independence restriction on unrestricted
directed edge effects. Cycle sums are the empirical content.

\item Kline's unrestricted edge model is not nested in the low-rank worker-firm wage
surface. AKM is nested as $\Lambda=0$, and a low-rank surface implies structured
edge effects, but arbitrary edge effects are a more general reduced form.
\end{enumerate}

\end{document}
